\documentclass[a4paper, serif, 10pt]{moderncv}

%% moderncv
% color and style
\moderncvcolor{black}
\moderncvstyle{banking}

%% page layout/margins
\usepackage[top=2.5cm, bottom=2.5cm, left=1.5cm, right=1.5cm]{geometry}

\nopagenumbers{}

%% Babel config for Italian language
% ref:
% - https://latex3.github.io/babel/guides/locale-italian.html
\usepackage[italian]{babel}

%% fontspec
\usepackage{fontspec}

\IfFontExistsTF{Linux Libertine}{
  \setmainfont[Ligatures={TeX, Common}, Numbers=OldStyle]{Linux Libertine}
}{}
\IfFontExistsTF{Linux Libertine O}{
  \setmainfont[Ligatures={TeX, Common}, Numbers=OldStyle]{Linux Libertine O}
}{}

%% CTeX
% CJK support, used to show CJK characters in the resume
%
% - heading=false: disable ctex heading
% - punct=kaiming: use kaiming punctuations styles for CJK
% - scheme=plain: use plain scheme, do not override \normalsize font size
% - space=auto: space settings for CJK characters
%
% CJK fonts are resolved per locale language, e.g. Noto Serif/Sans CJK JP
% for Japanese resumes, with fallback to the next region if not installed
%
% ref:
% - http://ctan.mirrorcatalogs.com/language/chinese/ctex/ctex.pdf
\usepackage[UTF8, heading=false, punct=kaiming, scheme=plain, space=auto]{ctex}

\IfFontExistsTF{Noto Serif CJK SC}{
  \setCJKmainfont{Noto Serif CJK SC}
}{}
\IfFontExistsTF{Noto Sans CJK SC}{
  \setCJKsansfont{Noto Sans CJK SC}
}{}

%% line spacing
\usepackage{setspace}
\setstretch{1.125}

%% URL styling - use normal text instead of monospace
\urlstyle{same}

%% Patch \httplink and \httpslink to handle full URLs
% The original moderncv macros always prepend http:// or https:// to URLs.
% This causes issues when the URL already has a protocol (e.g., https://example.com)
% resulting in malformed URLs like https://https://example.com.
% This patch checks if the URL already contains "://" and uses it directly if so.
% Note: We use \str_set:Nx (x-expansion) to expand macros like \@homepage first.
\ExplSyntaxOn
\str_new:N \g_yamlresume_url_str

\RenewDocumentCommand{\httpslink}{O{}m}{
  \str_set:Nx \g_yamlresume_url_str {#2}
  \str_if_in:NnTF \g_yamlresume_url_str {://}
    { \url{#2} }
    { \url{https://#2} }
}

\RenewDocumentCommand{\httplink}{O{}m}{
  \str_set:Nx \g_yamlresume_url_str {#2}
  \str_if_in:NnTF \g_yamlresume_url_str {://}
    { \url{#2} }
    { \url{http://#2} }
}
\ExplSyntaxOff

\name{Giulia Ferrari}{}
\title{Business Analyst Senior | Trasformazione Digitale}
\phone[mobile]{+39 02 9475 3281}
\email{giulia.ferrari@analisibusiness.it}
\homepage{https://www.giuliaferrari.it}

\address{Via Alessandro Volta 18, Milano, Lombardia, Italia, 20154}{}{}

\extrainfo{{\small \faLinkedin}\ \href{https://www.linkedin.com/in/giuliaferrari-ba}{@giuliaferrari-ba} {} {} {} • {} {} {} 
{\small \faGithub}\ \href{https://github.com/giuliaferrari-ba}{@giuliaferrari-ba}}

\begin{document}

\maketitle

\section{Informazioni personali}

\cvline{}{Business analyst con oltre 7 anni di esperienza nella raccolta requisiti, modellazione dei processi e delivery di soluzioni digitali in ambito bancario e retail.
%
\begin{itemize}
\item Esperta nella traduzione delle esigenze di business in user story, specifiche funzionali e flussi AS-IS/TO-BE, collaborando con stakeholder, product owner e team di sviluppo.
\item Solida conoscenza di metodologie Agile/Scrum, analisi dei dati, SQL e strumenti di visualizzazione come Power BI e Tableau.
\item Orientata ai risultati, con eccellenti capacità di comunicazione, facilitazione di workshop e gestione del cambiamento.
\end{itemize}}
\section{Formazione}

\cventry{set 2014–lug 2016}
        {Laurea magistrale, Economia e Management, Punteggio: 110/110 con lode}
        {Università Bocconi}
        {\url{https://www.unibocconi.it}}
        {}
        {\begin{itemize}
\item Percorso focalizzato sulla gestione dei processi, analisi dei dati e trasformazione digitale.
\item Tesi finale su modelli predittivi per l'ottimizzazione della customer experience nel retail.
\end{itemize}
\textbf{Corsi}: Analisi dei dati per il business, Sistemi informativi aziendali, Gestione dei processi aziendali, Ricerca operativa, Marketing analytics}

\cventry{set 2011–lug 2014}
        {Laurea triennale, Ingegneria Gestionale, Punteggio: 105/110}
        {Politecnico di Milano}
        {\url{https://www.polimi.it}}
        {}
        {\textbf{Corsi}: Economia aziendale, Statistica, Sistemi informativi, Ricerca operativa, Gestione della produzione}
\section{Esperienza professionale}

\cventry{gen 2021–Presente}
        {Senior Business Analyst}
        {Nexa Retail Italia}
        {\url{https://www.nexaretail.it}}
        {}
        {\begin{itemize}
\item Raccolta e prioritizzazione dei requisiti con stakeholder di business, operations e IT.
\item Modellazione dei processi AS-IS/TO-BE e definizione dei requisiti funzionali per piattaforme retail.
\item Coordinamento del backlog con product owner e team di sviluppo in ambiente Agile/Scrum.
\item Definizione di KPI e dashboard Power BI per il monitoraggio delle performance operative.
\end{itemize}
\textbf{Parole chiave}: Raccolta requisiti, Agile/Scrum, Modellazione processi, Power BI, Stakeholder management}

\cventry{mar 2018–dic 2020}
        {Business Analyst}
        {Banca Aurora}
        {\url{https://www.bancaaurora.it}}
        {}
        {\begin{itemize}
\item Analisi funzionale e redazione di specifiche per servizi di digital banking e portali self-service.
\item Scrittura di user story, criteri di accettazione e documentazione di supporto al rilascio.
\item Pianificazione ed esecuzione di test UAT con key user e supporto al team di sviluppo.
\item Monitoraggio dei difetti e coordinamento delle attività di remediation con il business.
\end{itemize}
\textbf{Parole chiave}: Analisi funzionale, SQL, User story, Test UAT, Digital banking}

\cventry{set 2016–feb 2018}
        {Junior Business Analyst}
        {Consulenza Sigma}
        {\url{https://www.consulenzasigma.it}}
        {}
        {\begin{itemize}
\item Supporto alla mappatura dei processi aziendali e alla produzione di documentazione funzionale.
\item Creazione di report e analisi Excel per il monitoraggio delle attività progettuali.
\item Partecipazione a workshop con il cliente per la raccolta e la validazione dei requisiti.
\end{itemize}
\textbf{Parole chiave}: Documentazione, BPMN, Excel avanzato, Reporting}
\section{Lingue}

\cvline{Italiano}{Madrelingua o bilingue \hfill \textbf{Parole chiave}: Madrelingua}
\cvline{Inglese}{Competenza professionale completa \hfill \textbf{Parole chiave}: IELTS 7.5, Business English}
\cvline{Francese}{Competenza professionale limitata \hfill \textbf{Parole chiave}: DELF B1}
\section{Competenze}

\cvline{Business Analysis}{Esperto \hfill \textbf{Parole chiave}: Raccolta requisiti, User story, BPMN, Analisi AS-IS/TO-BE, UAT}
\cvline{Metodologie}{Avanzato \hfill \textbf{Parole chiave}: Agile, Scrum, Kanban, Waterfall, Design thinking}
\cvline{Dati e Strumenti}{Avanzato \hfill \textbf{Parole chiave}: SQL, Power BI, Tableau, Jira, Confluence}
\cvline{Soft Skills}{Esperto \hfill \textbf{Parole chiave}: Comunicazione, Facilitazione, Negoziazione, Problem solving, Leadership}
\section{Premi}

\cventry{nov 2022}
        {Nexa Retail Italia}
        {Premio Innovazione Digitale}
        {}
        {}
        {Riconoscimento per il contributo alla riduzione del 20\% dei tempi di onboarding dei nuovi fornitori tramite la digitalizzazione dei processi.}

\cventry{giu 2019}
        {Banca Aurora}
        {Best Team Project}
        {}
        {}
        {Premio assegnato al team per il lancio di un portale self-service per oltre 150.000 clienti retail.}
\section{Certificazioni}

\cventry{mar 2023}
        {IIBA}
        {Certified Business Analysis Professional (CBAP)}
        {\url{https://www.iiba.org/certification/cbap/}}
        {}
        {}

\cventry{set 2021}
        {Scrum.org}
        {Professional Scrum Product Owner I}
        {\url{https://www.scrum.org/professional-scrum-product-owner-certifications}}
        {}
        {}

\cventry{mag 2022}
        {Microsoft}
        {Microsoft Certified: Power BI Data Analyst Associate}
        {\url{https://learn.microsoft.com/en-us/certifications/power-bi-data-analyst-associate/}}
        {}
        {}
\section{Pubblicazioni}

\cventry{ott 2021}
        {Modelli predittivi per la personalizzazione dell'offerta retail}
        {FrancoAngeli Editore}
        {\url{https://www.francoangeli.it/}}
        {}
        {\begin{itemize}
\item Articolo tecnico sulla segmentazione della clientela e sull'uso di modelli predittivi per migliorare la customer experience.
\item Analisi di casi di studio nel settore retail italiano.
\end{itemize}}
\section{Progetti}

\cventry{feb 2019–nov 2020}
        {Piattaforma web per la gestione in autonomia di operazioni bancarie quotidiane.}
        {Portale Self-Service Clienti}
        {\url{https://www.bancaaurora.it/self-service}}
        {}
        {\begin{itemize}
\item Definizione dei requisiti funzionali e dei flussi utente per oltre 30 funzionalità.
\item Coordinamento delle attività di UAT con 12 key user e supporto al team di sviluppo.
\item Riduzione del 25\% delle chiamate al contact center nei primi sei mesi post-lancio.
\end{itemize}
\textbf{Parole chiave}: Digital banking, User experience, UAT, Self-service}

\cventry{mar 2021–set 2021}
        {Sistema di reporting direzionale per monitorare vendite, margini e rotazione delle scorte.}
        {Dashboard Decisionale Retail}
        {\url{https://www.nexaretail.it/progetti/dashboard-retail}}
        {}
        {\begin{itemize}
\item Raccolta dei KPI con i responsabili di area e definizione del modello dati.
\item Realizzazione di dashboard Power BI con aggiornamento giornaliero e controlli di qualità.
\item Adozione da parte di 45 store manager e riduzione del tempo di reporting settimanale del 60\%.
\end{itemize}
\textbf{Parole chiave}: Power BI, KPI, Retail, Data visualization}
\section{Interessi}

\cvline{Sport}{Running, Yoga, Trekking}
\cvline{Cultura}{Lettura, Fotografia, Viaggi}
\section{Volontariato}

\cventry{gen 2022–Presente}
        {Mentor di analisi dati}
        {Associazione Digital Skills Italia}
        {\url{https://www.digitalskills.it}}
        {}
        {\begin{itemize}
\item Mentoraggio di studenti e giovani professionisti su analisi dei dati e strumenti di business intelligence.
\item Organizzazione di workshop su Power BI, SQL e storytelling con i dati.
\item Supporto alla definizione di percorsi formativi orientati alle esigenze del mercato.
\end{itemize}}
    
\end{document}